\section{模型优势与局限性}
\label{sec:strengths}

\subsection{优势}

\textbf{成分数据前置规范。} 三个子问题共用同一套近全零过滤标准与 CLR 变换，既处理了相对丰度的定和约束，也让三个子问题建立在同一特征标准上。这样疾病预测、标志物筛选和跨疾病分析之间可以直接衔接，不会因为预处理方式不同而难以比较。

\textbf{评估框架。} 考虑到 ACC 在类别不平衡的影响，模型评估以 AUC 为主指标，同时报告少数类 F1 与 Recall。性能用分层 5 折交叉验证估计，再用 LOOCV 做兜底，并以两者差距小于 0.025 判断是否存在过拟合。两套验证相互对照，避免小样本下把准确率误认为泛化能力。

\textbf{稳定性选择抗小样本方差。} Lasso 加 Bootstrap 稳定性选择用特征被选频率量化选择不确定性，比单次特征选择更稳健；Fisher 存在/缺失 和 Wilcoxon 非零丰度两路信号检验，配合 BH-FDR 校正控制假阳性。CRC 3/4、IBD 4/4 命中已知标志物，数据筛选结果有生物背景的研究文献背书。

\textbf{跨疾病评估与负结果归因。} 在 LODO 框架下比较四种跨疾病策略，逐一尝试多种模型后最优 AUC 0.6068 仍低于可用性参考线 0.65，本文如实报告负结果，并通过六轮归因实验把原因收敛到疾病特异信号与数据固有限制。

\subsection{局限性}

\textbf{小样本限制统计效力。} 单疾病样本量只有 110–253，过滤后特征仍有 264 维，$n\ll p$ 使标志物筛选和跨疾病结论的统计效力有限。虽然本文用稳定性选择、BH-FDR 和 LOOCV 做了保守处理，但结论仍需要更大规模独立数据验证。

\textbf{Obesity 弱信号。} Obesity 20 个稳定标志物中无 FDR 显著，可信度低，因此本文将其明确标注。

\textbf{批次效应。} 三数据集来自不同研究，虽 silhouette 0.070 近 0 不主导全局结构，但跨疾病性能仍受批次差异影响，主流程没有把 ComBat 作为正式预处理。但即便在归因实验中做过 ComBat 校正，结果显示校正后 AUC 反而下降，所以没有纳入主模型。

\textbf{第三问 跨疾病预测负结果。} 不同疾病的标志物高度特异，方向仅 51.2\% 一致，二项符号检验 p=0.5364 不显著；跨疾病预测模型需要更详细的判断标签。这个负结果本身反映了数据限制，也说明当前方法尚不能直接用于跨疾病预测。

\textbf{第二问 RF 佐证失效。} RF 的置换重要性在三组数据上的精度几乎全是 $10^{-17}$ 量级，无法有效排序，因此正文不提及RF的重要性结果。该现象的原因尚未查明，留待后续调查。
